\subsubsection{Elementos mecânicos}

A modelagem tridimensional da \textit{case} do dispositivo foi utilizado o software Autodesk Fusion, o qual possui a ferramenta de projeto assistido por computador (CAD, do inglês \textit{computer-aided design}, que oferece recursos de modelagem paramétrica, colaboração em nuvem e licenciamento gratuito para fins educacionais e de desenvolvimento. Para a fabricação do invólucro do protótipo foi adotada a manufatura aditiva pelo processo de Modelagem por Deposição Fundida (FDM, do inglês \textit{Fused Deposition Modeling}) devido ao seu baixo custo de implementação, rapidez na produção de peças e facilidade de prototipagem, características que favorecem o desenvolvimento iterativo do produto.

Para os materiais termoplásticos, foram utilizados dois materiais para etapas distintas do desenvolvimento. O Ácido Polilático (PLA, do inglês \textit{Polylactic Acid}) foi utilizado nas iterações iniciais do invólucro, em razão de sua facilidade de impressão, boa estabilidade dimensional e baixo custo

Para a versão final do protótipo, foi adotado o Polietileno Tereftalato Glicol (PETG, do inglês \textit{Polyethylene Terephthalate Glycol}). Essa é justificada por duas propriedades críticas para o contexto de aplicação do dispositivo: maior resistência térmica em relação ao PLA e tolerância à limpeza com álcool isopropílico, procedimento de higienização padrão em ambientes hospitalares e clínicos.

% CAD, justificada pela possibilidade de versionamento do projeto, pelo ambiente paramétrico que possibilita iteração rápida sobre dimensões, tolerâncias e encaixes e pela geração de documentação visual rica. 

% A modelagem tridimensional do invólucro do dispositivo foi realizada em software CAD, justificada por três razões principais. Primeiramente, a modelagem em CAD confere rastreabilidade ao projeto, uma vez que cada revisão do modelo é versionada e documentada, permitindo identificar com precisão quais alterações foram realizadas entre iterações e por qual motivo. Em segundo lugar, o ambiente paramétrico do CAD possibilita iteração rápida sobre dimensões, tolerâncias e encaixes, de modo que ajustes decorrentes de testes físicos podem ser incorporados ao modelo em tempo reduzido, sem necessidade de retrabalho completo da geometria. Por fim, a modelagem 3D gera documentação visual rica, que facilita a comunicação dos resultados tanto na seção de resultados deste trabalho.

% O software selecionado para a modelagem foi o Autodesk Fusion, escolha justificada pela integração nativa entre o ambiente CAD mecânico, pela disponibilidade de licenciamento gratuito para uso acadêmico e pelos recursos de colaboração em nuvem. Como alternativas de stack aberto, ferramentas como FreeCAD e OnShape poderiam ser adotadas em contextos onde a independência de fornecedor é um requisito, sem prejuízo significativo à funcionalidade necessária para este projeto.

% Para a fabricação do invólucro físico do protótipo, foi adotada a manufatura aditiva pelo processo de Modelagem por Deposição Fundida (FDM, do inglês \textit{Fused Deposition Modeling}). Essa tecnologia se mostra o caminho mais adequado para a etapa de prototipagem pelas seguintes razões: (i) baixo custo de produção por peça, uma vez que utiliza filamentos termoplásticos de ampla disponibilidade comercial e equipamentos acessíveis; (ii) prazo reduzido de fabricação, permitindo que uma nova iteração do casco seja produzida em poucas horas após a atualização do modelo CAD; e (iii) adequação às peças não estruturais do dispositivo, cujos requisitos mecânicos — alojamento de componentes eletrônicos, fixação por insertos e fechamento por parafusos — são plenamente atendidos pela resistência mecânica dos materiais impressos por FDM.

% Para os materiais termoplásticos, foram utilizados 2 materiais para etapas distintas do desenvolvimento. O Ácido Polilático (PLA, do inglês \textit{Polylactic Acid}) foi utilizado nas iterações iniciais do casco. Esse material apresenta facilidade de impressão, boa definição dimensional e custo reduzido, sendo adequado para validação rápida de forma, encaixes e ergonomia, sem exigência de propriedades físico-químicas avançadas. O Polietileno Tereftalato Glicol (PETG, do inglês \textit{Polyethylene Terephthalate Glycol}) foi adotado para a versão final do protótipo. A escolha do PETG é justificada por duas propriedades críticas para o contexto de aplicação do dispositivo: (i) maior resistência térmica em relação ao PLA, com temperatura de deflexão sob carga (HDT) superior, o que reduz o risco de deformação em ambientes com variação térmica; e (ii) tolerância à limpeza com álcool isopropílico (IPA), procedimento de higienização padrão em ambientes hospitalares e clínicos.

% Para a fixação e montagem do protótipo foi utilizado o Insertos rosqueados de latão M3, instalados a quente nas paredes do casco por meio de ferro de solda ou estação de retrabalho, garantindo roscas duráveis e resistentes ao torque de aperto repetido. A união dos componentes foi realizada por meio de parafusos de cabeça cilíndrica com sextavado interno (do inglês \textit{socket head cap screw}), em aço inoxidável, nos comprimentos de 6~mm, para a fixação de tampa, e 10~mm, para a fixação de componentes internos ao SBC e ao display. As Porcas M3 sextavadas foram utilizadas em pontos onde a inserção de insertos não é viável geometricamente, alojadas em cavidades hexagonais impressas no próprio casco. Por fim, os Espaçadores M3 em latão, do tipo fêmea-fêmea nos comprimentos de 5~mm e 10~mm, foram utilizados para elevação e isolamento mecânico da placa SBC e da placa de interface em relação ao fundo do casco.


% % Justificar a modelagem 3D em CAD pela rastreabilidade do projeto, possibilidade de iteração rápida e documentação visual para a banca. Expor que a manufatura aditiva por FDM é o melhor caminho para o protótipo pelo baixo custo, prazo curto e adequação a peças não estruturais.
% % Software CAD: Fusion 360 (alternativas: OnShape, FreeCAD para defesa de stack open-source).
% % Material: PLA para iteração e PETG para a versão final do protótipo — justificar PETG pela maior resistência térmica e tolerância à limpeza com álcool isopropílico (ambiente hospitalar).
% % Listar e especificar:

% % Insertos rosqueados de latão (M3 padrão)
% % Parafusos, porcas, espaçadores (cabeçote, comprimento)
% % Suportes internos para fixação do SBC, sensor, display e bateria
% % Cases base inferior (eletrônica) e superior (display + acesso ao sensor)

% % Detalhar processo de montagem com sequência fotográfica para a seção de resultados. Documentar pelo menos duas iterações de casco evidenciando refinamentos.